% Desenvolvido por André Leite - DE - UFPE
% Ajustes Raydonal Ospina - DE - UFPE/UFBA

\documentclass[a4paper,11pt,oneside,twocolumn]{Config/milktest}

%---------------------- customização ----------------------



\usepackage{fontspec}
\usepackage{xunicode}
\usepackage{xltxtra}
\usepackage[LGR,T1]{fontenc}
\usepackage{ebgaramond}

%\usepackage[cp1252]{inputenc}                %% Use um dos dois esquemas
\usepackage[english,brazil]{babel}                 %% Hifenizacao e estilo
\usepackage{pstricks}                          %% Pacote de comandos p/ Facilitar o Postscripts
\usepackage{pstricks-add}                    %% " "
\usepackage{pst-all,pst-math,pst-plot, pst-func}       %% " "
\usepackage{textcomp}							%%
%\usepackage[tight,nice]{units}				%%
\usepackage[centertags]{amsmath}       %% [leqno]
\usepackage{amsfonts}                         %% Pacotes da AMS
\usepackage{amssymb}                         %%
\usepackage{stmaryrd} 							%%
\usepackage{Config/milkstyle}
%\usepackage{sabon}

\usepackage{nicefrac}%  Nice Frescura
\usepackage{balance}%
\usepackage{lastpage} %                      %% Referênciar a última página
\usepackage{graphicx}                        %% Mais ferramentas para
%\usepackage[T1]{pbsi} % Procurar e instalar no Mac!!!!! <------<-------<------
\usepackage{mathrsfs}                        %% $\mathscr{ABCD}$
\usepackage{lettrine}                        %% ...
%\usepackage[normal,bf,up]{caption2}          %% Controlar melhor os captions....
\usepackage{eucal}                           %% beau
\usepackage{colortbl} %
\usepackage{hyperref}                 %%
%\usepackage{amsthm}
\usepackage{indentfirst}%
\usepackage{enumitem}%
\usepackage{colordvi}                        %% Usar Cores...
\usepackage{fancyhdr}%
\usepackage{yfonts}%
\usepackage{theorem}%
\usepackage{multirow}                        %% Juntar linhas em tabelas
\usepackage{midfloat}%
\usepackage{cancel}%
%\usepackage[square]{natbib}%sectionbib,
\usepackage{fancyhdr} % cabeçalhos ...
\usepackage{rotating}


\usepackage{lipsum} 
%% Lista de Pacotes utilizados


%% Deixando o caption mais interessante...
%\renewcommand{\captionfont}{\itshape}


\hypersetup{%backref,%
%            pdfpagemode=FullScreen,%
            colorlinks=true,%
            citecolor=blue,%
            urlcolor=blue, %
            % urlcolor=black, %
            linkcolor=blue,%
            pdfauthor={André Leite Wanderley}, %
            pdftitle={1EE/ET406}, %
            pdfcreator={André Leite Wanderley}, %
            breaklinks=true, %
            linktocpage=true
            }%


\DeclareMathOperator*{\Inf}{Inf}%
\DeclareMathOperator*{\Sup}{Sup}%
\DeclareMathOperator*{\Min}{Min}%
\DeclareMathOperator*{\Max}{Max}
\DeclareMathOperator*{\Sgn}{Sgn}%
\DeclareMathOperator*{\E}{E}%
\newcommand{\mbf}{\boldsymbol}%
\DeclareMathAlphabet\mathbb{U}{fplmbb}{m}{n}%
%
%\numberwithin{equation}{section}                       %%
%\numberwithin{table}{section}                          %% Numeração dos floats...
%\numberwithin{figure}{section}                         %%


\newcommand{\oo}{\mathbin{\bullet}}
\newcommand{\oc}{\mathbin{\circ}}
\newcommand{\noo}{\mathbin{\not\!\bullet}}
\newcommand{\cf}{\emph{cf.}\ }
\newcommand{\ie}{\emph{i.e.}}
\newcommand{\eg}{\emph{e.g.}}
\newcommand{\raab}{\emph{reductio ad absurdum}}
\newcommand{\eetc}{\textit{et c\ae tera}}
\newcommand{\dom}{\mathfrak{D}}
\newcommand{\ima}{\mathfrak{I}}
\newcommand{\idenA}{\mathcal{I}_\mathcal{A}}
\newcommand{\ce}{\ensuremath{\mathcal{X}}}
\newcommand{\piq}{\ensuremath{\succsim}}
\newcommand{\pq}{\ensuremath{\succ}}
\newcommand{\iq}{\ensuremath{\sim}}
\newcommand{\niq}{\ensuremath{\nsim}}
\newcommand{\npiq}{\ensuremath{\not\succsim}}
\newcommand{\npq}{\ensuremath{\nsucc}}
\newcommand{\po}[2]{\ensuremath{(#1,#2)}}

%\setlist{noitemsep}
%\setlist{labelindent=\parindent} % << Usually a good idea
\setlist{itemsep=0pt}%
\setitemize{leftmargin=\parindent} %
\setitemize[1]{label=$\triangleright$} %
%\setenumerate[1]{label=\arabic*.,ref=\arabic*}
\setenumerate[1]{label=\arabic*.,ref=\arabic*}
\setenumerate[2]{label=\theenumi.\arabic*.,ref=\theenumi.\arabic*}
\setenumerate[3]{label=\roman*.,ref=\theenumii.\roman*}
%\setdescription{%font=\sffamily\bfseries}

%                }%


%%%%%%%%%%%%%%%%%%%%%%%%%%%%%%%%
% Para facilitar minha vida

\newcommand{\Ds}{\displaystyle}%
\newcommand{\Ts}{\textstyle}%
\newcommand{\Ss}{\scriptstyle}%
\newcommand{\SSs}{\scriptscriptstyle}%

\def\aureos{%
            \newrgbcolor{gray1}{.9 .9 .9}
            \newrgbcolor{guess1}{1 0.388235294 0.090196078}
            \psset{unit=1in,cornersize=absolute}%
            \psline[linecolor=guess1,linewidth=1.2pt](-0.125,-1.291667)(1.1125,-1.291667)
            \psline[linecolor=guess1,linewidth=0.6pt](0.347683,-1.291667)(0.347683,-0.52685)
            \psline[linecolor=guess1,linewidth=0.6pt](0.347683,-0.999533)(-0.125,-0.999533)
            \psline[linecolor=guess1,linewidth=0.6pt](0.167134,-0.999533)(0.167134,-1.291667)
            \psline[linecolor=guess1,linewidth=0.6pt](0.167134,-1.111118)(0.347683,-1.111118)
            \psline[linecolor=guess1,linewidth=0.6pt](0.236098,-1.111118)(0.236098,-0.999533)
            \psline[linecolor=guess1,linewidth=0.6pt](0.236098,-1.068496)(0.167134,-1.068496)
            \psline[linecolor=guess1,linewidth=0.6pt](0.209756,-1.068496)(0.209756,-1.111118)
            \psline[linecolor=guess1,linewidth=0.6pt](0.209756,-1.084776)(0.236098,-1.084776)
           }%

\def\GuEsS{%
            \newrgbcolor{guess1}{1 0.388235294 0.090196078}
            \newrgbcolor{guess2}{0.960784314 0.109803922 0.043137255}
            \psset{unit=1in,cornersize=absolute}%
            \rput(0,-1.5){{\fontsize{35}{0}\usefont{OT1}{cmr}{m}{n}\selectfont\textsc{\textbf{S}}}}%
            \rput(0.625,-1.5){\shortstack{{\huge\textbf{\textsc{ystems}}}\\
                {\scriptsize\sc\hfill Research Group}}}
            \psline[linecolor=black,fillcolor=guess2,fillstyle=none,linewidth=3.6pt](-0.125,-1.725)(1.1125,-1.725)%
            \uput{0.5ex}[l](1.148214,-1.229167){\llap{{\tiny\sc UFBA}}}%
            \psline[linecolor=guess1,fillcolor=guess1,fillstyle=none,linewidth=1.2pt](-0.125,-1.291667)(1.1125,-1.291667)%
            \psline[linecolor=guess1,fillcolor=guess1,fillstyle=none,linewidth=.6pt](0.347683,-1.291667)(0.347683,-0.52685)%
            \psline[linecolor=guess1,fillcolor=guess1,fillstyle=none,linewidth=.6pt](0.347683,-0.999533)(-0.125,-0.999533)%
            \psline[linecolor=guess1,fillcolor=guess1,fillstyle=none,linewidth=.6pt](0.167134,-0.999533)(0.167134,-1.291667)%
            \psline[linecolor=guess1,fillcolor=guess1,fillstyle=none,linewidth=.6pt](0.167134,-1.111118)(0.347683,-1.111118)%
            \psline[linecolor=guess1,fillcolor=guess1,fillstyle=none,linewidth=.6pt](0.236098,-1.111118)(0.236098,-0.999533)%
            \psline[linecolor=guess1,fillcolor=guess1,fillstyle=none,linewidth=.6pt](0.236098,-1.068496)(0.167134,-1.068496)%
            \psline[linecolor=guess1,fillcolor=guess1,fillstyle=none,linewidth=.6pt](0.209756,-1.068496)(0.209756,-1.111118)%
            \psline[linecolor=guess1,fillcolor=guess1,fillstyle=none,linewidth=.6pt](0.209756,-1.084776)(0.236098,-1.084776)%
           }%

%%%%%%%%%%%%%%%%%%%%%%%%%%%%%%%%%%%%%%%%%%%%%%%%%%%%%%%%


\theoremstyle{plain}                     %%
%\theoremstyle{margin}                   %%
\theoremheaderfont{\scshape\bfseries}    %%
\theorembodyfont{\rmfamily}                       %%

\newtheorem{Prop}{\small Proposition}[section]  %%
\newtheorem{Def}{\small Definition}[section]    %%
\newtheorem{Not}{\small Notice}[section]    %%
\newtheorem{Theo}{\small Theorem}[section]      %%
\newtheorem{Lem}{\small Lema}[section]         %%
\newtheorem{Coro}{\small Corollary}[section]    %%
\newtheorem{Conj}{\small Conjecture}[section]   %%
\newtheorem{Prob}{\small Problem}[section]      %%
\newtheorem{Exa}{\small Example}[section]       %%
%\theoremheaderfont{\bfseries}    %%
%\theorembodyfont{\rmfamily}  %%
\newtheorem{proof}{\small Demonstration}[section]
\newtheorem{sol}{\small Solution}[section]


\newenvironment{Proofm}               % Deve ter uma maneira
    {%\strut \par\medskip             % ridiculamente mais
     %\strut\hspace{\parindent}      % fácil de fazer isso...
     %\begin{minipage}{.8\linewidth} % I'm sure about that.
     \begin{proof}\small}                   %
    {\strut\hfill\textsc{q.e.d.}\end{proof}                      %
    %\end{minipage}                  %
    %\strut\par\medskip
                        }             %

\newenvironment{solm}               % Deve ter uma maneira
    {%\strut\par%\medskip             % ridiculamente mais
     %\strut\hspace{\parindent}      % fácil de fazer isso...
     %\begin{minipage}{.8\linewidth} % I'm sure about that.
     %\begin{sol}
     \small
     \noindent\textbf{Solução:}
     %\par %\smallskip
     }                   %
    {\strut\hfill\textsc{$\blacksquare$}%\end{sol}                      %
    %\end{minipage}                  %
    %\strut\par\medskip
                        }             %



\newenvironment{Propm}               % Deve ter uma maneira
    {\strut\par\smallskip             % ridiculamente mais
     \strut\hspace{\parindent}      % fácil de fazer isso...
     \begin{minipage}{.8\linewidth} % I'm sure about that.
     \begin{Prop}\small}                   %
    {\end{Prop}                      %
    \end{minipage}                  %
    \strut\par\medskip
    }             %


\newenvironment{Defm}               % Deve ter uma maneira
    {\strut\par\smallskip             % ridiculamente mais
     \strut\hspace{\parindent}      % fácil de fazer isso...
     \begin{minipage}{.8\linewidth} % I'm sure about that.
     \begin{Def}\small}                   %
    {\end{Def}                      %
    \end{minipage}                  %
    \strut\par\medskip}             %

\newenvironment{Theom}               % Deve ter uma maneira
    {\strut\par\smallskip             % ridiculamente mais
     \strut\hspace{\parindent}      % fácil de fazer isso...
     \begin{minipage}{.8\linewidth} % I'm sure about that.
     \begin{Theo}\small}                   %
    {\end{Theo}                      %
    \end{minipage}                  %
    \strut\par\medskip}

\newenvironment{Notm}               % Deve ter uma maneira
    {\strut\par\smallskip             % ridiculamente mais
     \strut\hspace{\parindent}      % fácil de fazer isso...
     \begin{minipage}{.8\linewidth} % I'm sure about that.
     \begin{Not}\small}                   %
    {\end{Not}                      %
    \end{minipage}                  %
    \strut\par\medskip}


\newenvironment{Exam}               % Deve ter uma maneira
    {%\strut\par\smallskip             % ridiculamente mais
     %\strut\hspace{\parindent}      % fácil de fazer isso...
     %\begin{minipage}{1\linewidth} % I'm sure about that.
     \begin{Exa}\small}                   %
    { \scriptsize$\blacksquare$\end{Exa}                      %
    %\end{minipage}                  %
    \strut\par\medskip}

\def\Pt#1(#2,#3){
       \expandafter\def\csname #1\endcsname{#2,#3}%
       \expandafter\def\csname #1X\endcsname{#2}%
       \expandafter\def\csname #1Y\endcsname{#3}}%

\newcommand{\hlinem}[3]{\noalign{\vspace{#3}\hrule height #1 depth #2}}
\newcommand{\hlinemilk}[3]{\noalign{}}

\newcommand{\quem}[1]{\par\medskip\strut\hfill\bf  --- #1}
\newcounter{questionNumber}
\newcommand{\question}{\stepcounter{questionNumber}\lettrine[lhang=.75, findent=10pt, nindent=0pt,lines=3, lraise=.2]{\thequestionNumber}{$\blacktriangleright$}}


\newenvironment{milkquote}
    {\medskip %
    \begin{center}
    \begin{minipage}{.9\linewidth}
    \hspace{.25cm}\small\itshape}%
    {\medskip\end{minipage}
    \end{center}
    }%


\newcommand{\fim}{%
                 \par\medskip
                 \begin{center}
                      \includegraphics[width=1cm]{ornaments/ornamento4.eps}
                 \end{center}
                 \vspace*{.5cm}
                 }%



\newenvironment{milkanswer}
    {\par %
    \begin{center}
    \color{orange}\textbullet\quad\textbf{Solution}\quad\textbullet
    \end{center}\small
    }%
    {
    \fim
    }


%%%
%%% Set up page layout parameters.
%%%
\setlength\paperheight{297mm}      %
\setlength\paperwidth{210mm}       % papel a4
%\setlength{\parindent}{6mm}

\addtolength{\hoffset}{-1in}%
\addtolength{\voffset}{-1.5in}

\setlength{\oddsidemargin}{25mm}    % 
\setlength{\topmargin}{5cm}        %
\setlength{\textwidth}{170mm}      % Margem direita: 25mm
\setlength{\textheight}{220mm}  % Margem Inferior: 25mm
\setlength{\marginparsep}{0mm}     %

\setlength{\headheight}{16mm}
\renewcommand{\columnsep}{1cm}%
%\marginsize{1.25in}{1.75in}{1.25in}{1.25in}



% If you can to leave the header here
%---

%---

\renewcommand{\headrulewidth}{0pt}
\definecolor{cor}{rgb}{.5,0, 0}

%%%%%%%%%%%%%%%%%%%%%%%%%%% Font 

\defaultfontfeatures{Ligatures=Historic,Contextuals=Alternate, Numbers=OldStyle,RawFeature={ +hlig, +dlig, +ss05}}
\setmainfont[RawFeature={+ss05, +hlig, +dlig}, BoldFont = Gentium Basic Bold, ItalicFeatures={RawFeature={+clig},CharacterVariant=5:0}]{EB Garamond}
% % Do the replacements manually in titles, to not ruin the textsc
% \newfontfamily\booktitlefont[RawFeature={-ss02},LetterSpace=40,WordSpace=6]{EB Garamond}
% \newfontfamily\spacedfont[RawFeature={-ss02},LetterSpace=20,WordSpace=3]{EB Garamond}
% \newfontfamily\lettrinefont{EB Garamond Initials}
% \newfontfamily\headerfont{EB Garamond}

%\setmainfont[BoldFont = Gentium Basic Bold,
%      ItalicFont = Gentium Plus Italic,
%      BoldItalicFont = Gentium Basic Bold Italic,
%      Mapping = tex-text]{Gentium Plus}





 

%%%%%%%%%%%%%----- Header------%%%%%%%%%%%%%%%%%%%
\newsavebox{\mygraphic}
\sbox{\mygraphic}{\includegraphics[scale=1.2]{Figuras/ufba.eps}}

\fancypagestyle{milktablestyle}{
    \fancyhead{}
    \fancyfoot[C]{\thepage/\pageref{LastPage}} %
}

\pagestyle{fancy}
\fancyhead{}
\fancyhead[C]{
\rput(0,2){
  \begin{minipage}{1.75cm}\centering
    \includegraphics[width=1.75cm]{Figuras/ufbaBlack.eps}
  \end{minipage}~\hfill~\begin{minipage}{10cm}\Large\centering
    \color{cor}{\scshape universidade federal da bahia}\\
    {\scshape instituto de matemática e estatística}\\
    {\scshape departamento de estatística}\\[10pt]
    % Aqui as informações da prova e a data
    {\Large\scshape  prova de 2da chamada --- mat 236\\ {\large 14 de dezembro de 2022}} 
  \end{minipage}~\hfill~\begin{minipage}{1.75cm}\centering
    \includegraphics[width=2.0cm]{Figuras/ufba.eps}
  \end{minipage}
  }}

\fancyfoot[C]{\thepage/\pageref{LastPage}} % 

%%%%%%%%%%%%%----- Header------%%%%%%%%%%%%%%%%%%%


\begin{document}



% %------------------------------------------------
\beb
{
\begin{description}
\item[Professor:] Raydonal Ospina
\item[E-mail:] raydonalmartinez@ufba.br
\item[Site do curso:] \url{https://castlab.org/courses/mat236/}

\item[Regras:] Leia com atenção as perguntas. Todas as questões devem ser DETALHADAS. A prova deve ser claramente resolvida. Seja claro e organizado. Pode fazer uso da calculadora. É expressamente {\bf proibido} o uso do celular durante a prova.

% \item[Atenção:] Responsa única e exclusivamente o número de questões {\bf suficientes} para completar 10 (dez) pontos.  Fique a vontade para escolher as perguntas que achar mais interessantes. Especifique sua escolha, por exemplo, 2(a). (isto quer dizer que você escolheu  o numeral (a) da questão 2).
\end{description}
}
\eeb

% %%%%%%%%%%%%%%%%%%%%%%%%%%%%%%%%%%%%%%%%%%%%%%%%%%%%%%%%%%%%%%%%%
\balance
%%%%%%%%%%%%%%%%%%%%%%%%%%%%%%%%%%%%%%%%%%%%%%%%%%%%%%%%%%%%%%%%%

\medskip 
\question~\textbf{[2pts]}  Estuda-se a relação linear entre duas variáveis observadas $x$ e $y$. Uma amostra de $20$ pares dessas variáveis forneceu os seguintes valores:
$$
\sum_{i=1}^{20} x_{i}= 600,\quad \sum_{i=1}^{20} y_{i}= 2150, 
$$ 
$$
\sum_{i=1}^{20} x^{2}_{i}= 18.662,\quad \sum_{i=1}^{20} y^{2}_{i}= 235.270,
$$ 
$$
\sum_{i=1}^{20} x_{i}y_{i}= 65.921
$$
Determine a correlação e ajuste uma reta aos dados.


% \medskip 
% \question~\textbf{[4pts]} Sobre o modelo de regressão linear simples, verifique se as alternativas a seguir são corretas ou falsas:
% \begin{description}
% \item [\ (a)]  Os estimadores do modelo de regressão linear simples minimizam a soma dos resíduos.
% \item[ \ (b)] Um elevado coeficiente de determinação referente a um modelo de regressão linear simples para uma amostra não muito grande, não implica, necessariamente, que a reta de regressão passe próxima a todos os pontos amostrados e que o modelo esteja bem ajustado
% \item[ \ (c)] Em um modelo linear simples a variância $\sigma^2$ e o valor esperado dos erros devem ser zero?
% \item[ \ (d)] O modelo de regressão é ``linear'' porque é linear a relação entre o parâmetro de inclinação da variável independente e a variável resposta.
% \end{description}




% \medskip 
% \question~\textbf{[3pts]} O resíduo em um modelo de regressão linear simples é definido como $r_i=y_i-\widehat{y_i},$ em que $\widehat{y_i}=\widehat{\beta_0}+\widehat{\beta_1}x_i,$ $i=1,\ldots,n$ 
% \begin{description}
%     \item [ \ (a)] $\sum_{i=1}^n r_i =0$
%     \item [ \ (b)] $\sum_{i=1}^n y_i = \sum_{i=1}^n {\widehat y_i}$
%     \item [ \ (c)] O resíduos são sempre independentes.
% \end{description}



% \medskip 
% \question~\textbf{[4pts]} Consideremos a avaliação de associação linear entre  o tempo de reação $y_i$ do $i$-ésimo indivíduo e $x_i$ a sua respectiva idade. O ajuste do modelo aos dados forneceu a seguinte equação da reta estimada 
% $${\widehat y}_i = 0,50 + 0,9 0x_i
% \quad i=1,\ldots.n$$

% Adicionalmente, a análise de variância (ANOVA) do  modelo de regressão ajustado pode ser resumida no seguinte quadro

% \begin{center}
% \begin{tabular}{cccccc} \hline
% F. V.       & g.l. & S.Q.     & Q.M.    & F       \\ \hline
% Regressão   & 1    & 810    & 810   & 25,90   \\ \hline
% Resíduo     & 18   & 563    & 31,28 &         \\ \hline
% Total       & 19   & 1.373  & 72,26 &         \\ \hline
% \end{tabular}
% \end{center}
% em que $F\sim {\cal F}_{(1, n-2)}$ e para um nivel de significancia $\alpha=5\%$ temos que o quantil da distribuição  ${\cal F}(0.05)_{(1, 18)}$ é 4.41.
% \begin{description}
% \item [ \ (a)] Você acha que conhecer a idade do individuo ajuda a prever o tempo de reação?
% \item [ \ (b)] Você acha que a redução nos resíduos foi grande?
% \item [ \ (c)] Qual o valor do coeficiente de determinação $R^2$ ? Interprete esse número.
% \item [ \ (d)]  Se fosse concluído que podemos considerar a inclinação $\beta_1= 0$, como deveria ser reescrito o
% modelo ajustado? Justifique.
% \end{description}

\medskip 
\question~\textbf{[2pts]} Uma companhia aérea sabe que 5\% das
pessoas que fizeram reserva para um voo
não apareceram para viagem. Consequentemente, a política da empresa é vender 52
bilhetes para um voo que pode suportar 50
passageiros. Qual é a probabilidade de que
exista uma poltrona para cada passageiro
que aparece?

\medskip 
\question~\textbf{[4pts]}
% ~Defina, \textit{valor-p} ou nível descritivo, Erro Tipo I e Erro Tipo II. Explique, com suas palavras,  as etapas para a realização de um \textit{teste de hipóteses}.
Considere uma amostra aleatória de 25 estudantes, aos quais foi perguntado o número de filmes assistidos  na semana anterior. O resultado é apresentado na Tabela~\ref{tab:ex1amostra}.

\begin{table}[!h]
\begin{center}\small\sffamily
		  \begin{tabular}{c|c|c|c}
           	     \hlinem{.5pt}{.5pt}{0mm}
		  N. de filmes  &  $f$ & $f_{\text{rel}}$ & $F_{\text{rel}}$\\
		    \hlinem{.5pt}{1pt}{0mm}
		    $0$ & $5$ & &\\		    
		    $1$ & $9$ & &\\		    
		    $2$ & $6$ & &\\
		    $3$ & $4$  & &\\  		    
		    $4$ & $1$  & &\\  		    
           	     \hlinem{.5pt}{.5pt}{0mm}
		  \end{tabular}
\caption{Resultado da pesquisa.}
\label{tab:ex1amostra}
\end{center}
\end{table}

Note-se que na Tabela~\ref{tab:ex1amostra}  denota-se por $f$  a frequência, por $f_{\text{rel}}$ a frequência relativa e $F_{\text{rel}}$ a frequência relativa acumulada. 
\begin{description}
\item[(a)]  Complete as colunas da Tabela.
\item[(b)]  Encontre o primeiro quartil.
\item[(c)] Encontre a mediana.
\item[(d)] Que porcentagem dos estudantes viram menos de 3 filmes?
\end{description}


\medskip 
\question~\textbf{[2pts]} 
O ajuste de um modelo de regressão linear simples produz os seguintes gráficos de diagnóstico (Figuras~\ref{f1} e \ref{f2}) que podem ser usadas para avaliar a adequabilidade do modelo.

\begin{figure}[!h]
\caption{\label{f1}Gráfico de valores ajustados contra os resíduos}
\centering\includegraphics[width=5.0cm]{RV.png}
\end{figure}


\begin{figure}[!h]
\caption{\label{f2}Gráfico de quantis normais}
\centering\includegraphics[width=5.0cm]{Rq.png}
\end{figure}

O que você pode concluir sobre o ajuste do modelo pela obervação destes gráficos? Justifique

% % Formulas
% \medskip 
% \section*{Fórmulas}
% \begin{align*}
% \hat{\beta_1}&=\frac{\sum\limits_{i=1}^{n} (x_i -\overline{x})(y_i - \overline{y})}{\sum\limits_{i=1}^{n} (x_i -\overline{x})^{2}} \\
% \hat{\beta_0}&=\overline{y}-\hat{\beta_1}\overline{x}  \\
% \end{align*}

% Tables

% \input{Tables/normal.tex}

% \newpage
\pagebreak

%\thispagestyle{empty}
\thispagestyle{milktablestyle}

\vspace*{-10mm}
\rput(3,3){
\psscalebox{.75}{
\psset{xunit=1.25cm,yunit=10cm,labelFontSize=\scriptsize}% 
\begin{pspicture}(0,0)(4,.25) 
\def\vari{.35}
\pscustom[fillstyle=vlines,hatchwidth=0.1pt,hatchsep=1pt,linestyle=none]{%
\psTDist[linecolor=black]{1.5}{4}    
\psline(4,0)(1.5,0)
}%   
%\psaxes[labels=none,ticks=none,% 
%	arrowscale=1.2]{-}(0,0)(4,0)%
\psaxes[labels=none,ticks=none, arrowscale=1.2]{-}(0,0)(-4.25,0)(4.25,0.0)
\psTDist[linecolor=green!60!black,linewidth=1pt,plotpoints=100]{-4}{4}    
\rput[c]{0}(0,-.03){\scriptsize\Red{$0$}}% 
\rput[c]{0}(1.5,-.03){$t^*$}% 
\psline[linecolor=red](0,0)(0,-.01)% 
\psline(1.5,0)(1.5,-.01)% 
\rput[l]{0}(2.5,0.2){\parbox{6.5cm}{Probabilidade $p$}}%
\psline{o-}(2,.03)(2.5,.17)
\end{pspicture}}%
}

\begin{table}[ht]
\vspace*{-10mm}
\centering
\caption{Valores críticos da distribuição $t$.}
\begin{tabular}{r||rrrrrrrrrrrr}
\hline
 \multirow{2}{*}{gl} & \multicolumn{9}{|c}{Probabilidade de cauda superior $p$} \\
  \cline{2-10}
 & 0.25 & 0.20 & 0.15 & 0.10 & 0.05 & 0.025 & 0.02 & 0.01 & 0.005 \\ 
\hline
1 & 1.0000 & 1.3764 & 1.9626 & 3.0777 & 6.3138 & 12.706 & 15.895 & 31.821 & 63.6567 \\ 
  2 & 0.8165 & 1.0607 & 1.3862 & 1.8856 & 2.9200 & 4.3027 & 4.8487 & 6.9646 & 9.9248 \\ 
  3 & 0.7649 & 0.9785 & 1.2498 & 1.6377 & 2.3534 & 3.1824 & 3.4819 & 4.5407 & 5.8409 \\ 
  4 & 0.7407 & 0.9410 & 1.1896 & 1.5332 & 2.1318 & 2.7764 & 2.9985 & 3.7469 & 4.6041 \\ 
  5 & 0.7267 & 0.9195 & 1.1558 & 1.4759 & 2.0150 & 2.5706 & 2.7565 & 3.3649 & 4.0321 \\ 
  6 & 0.7176 & 0.9057 & 1.1342 & 1.4398 & 1.9432 & 2.4469 & 2.6122 & 3.1427 & 3.7074 \\ 
  7 & 0.7111 & 0.8960 & 1.1192 & 1.4149 & 1.8946 & 2.3646 & 2.5168 & 2.9980 & 3.4995 \\ 
  8 & 0.7064 & 0.8889 & 1.1081 & 1.3968 & 1.8595 & 2.3060 & 2.4490 & 2.8965 & 3.3554 \\ 
  9 & 0.7027 & 0.8834 & 1.0997 & 1.3830 & 1.8331 & 2.2622 & 2.3984 & 2.8214 & 3.2498 \\ 
  10 & 0.6998 & 0.8791 & 1.0931 & 1.3722 & 1.8125 & 2.2281 & 2.3593 & 2.7638 & 3.1693 \\ 
  11 & 0.6974 & 0.8755 & 1.0877 & 1.3634 & 1.7959 & 2.2010 & 2.3281 & 2.7181 & 3.1058 \\ 
  12 & 0.6955 & 0.8726 & 1.0832 & 1.3562 & 1.7823 & 2.1788 & 2.3027 & 2.6810 & 3.0545 \\ 
  13 & 0.6938 & 0.8702 & 1.0795 & 1.3502 & 1.7709 & 2.1604 & 2.2816 & 2.6503 & 3.0123 \\ 
  14 & 0.6924 & 0.8681 & 1.0763 & 1.3450 & 1.7613 & 2.1448 & 2.2638 & 2.6245 & 2.9768 \\ 
  15 & 0.6912 & 0.8662 & 1.0735 & 1.3406 & 1.7531 & 2.1314 & 2.2485 & 2.6025 & 2.9467 \\ 
  16 & 0.6901 & 0.8647 & 1.0711 & 1.3368 & 1.7459 & 2.1199 & 2.2354 & 2.5835 & 2.9208 \\ 
  17 & 0.6892 & 0.8633 & 1.0690 & 1.3334 & 1.7396 & 2.1098 & 2.2238 & 2.5669 & 2.8982 \\ 
  18 & 0.6884 & 0.8620 & 1.0672 & 1.3304 & 1.7341 & 2.1009 & 2.2137 & 2.5524 & 2.8784 \\ 
  19 & 0.6876 & 0.8610 & 1.0655 & 1.3277 & 1.7291 & 2.0930 & 2.2047 & 2.5395 & 2.8609 \\ 
  20 & 0.6870 & 0.8600 & 1.0640 & 1.3253 & 1.7247 & 2.0860 & 2.1967 & 2.5280 & 2.8453 \\ 
  21 & 0.6864 & 0.8591 & 1.0627 & 1.3232 & 1.7207 & 2.0796 & 2.1894 & 2.5176 & 2.8314 \\ 
  22 & 0.6858 & 0.8583 & 1.0614 & 1.3212 & 1.7171 & 2.0739 & 2.1829 & 2.5083 & 2.8188 \\ 
  23 & 0.6853 & 0.8575 & 1.0603 & 1.3195 & 1.7139 & 2.0687 & 2.1770 & 2.4999 & 2.8073 \\ 
  24 & 0.6848 & 0.8569 & 1.0593 & 1.3178 & 1.7109 & 2.0639 & 2.1715 & 2.4922 & 2.7969 \\ 
  25 & 0.6844 & 0.8562 & 1.0584 & 1.3163 & 1.7081 & 2.0595 & 2.1666 & 2.4851 & 2.7874 \\ 
  26 & 0.6840 & 0.8557 & 1.0575 & 1.3150 & 1.7056 & 2.0555 & 2.1620 & 2.4786 & 2.7787 \\ 
  27 & 0.6837 & 0.8551 & 1.0567 & 1.3137 & 1.7033 & 2.0518 & 2.1578 & 2.4727 & 2.7707 \\ 
  28 & 0.6834 & 0.8546 & 1.0560 & 1.3125 & 1.7011 & 2.0484 & 2.1539 & 2.4671 & 2.7633 \\ 
  29 & 0.6830 & 0.8542 & 1.0553 & 1.3114 & 1.6991 & 2.0452 & 2.1503 & 2.4620 & 2.7564 \\ 
  30 & 0.6828 & 0.8538 & 1.0547 & 1.3104 & 1.6973 & 2.0423 & 2.1470 & 2.4573 & 2.7500 \\ 
  40 & 0.6807 & 0.8507 & 1.0500 & 1.3031 & 1.6839 & 2.0211 & 2.1229 & 2.4233 & 2.7045 \\ 
  50 & 0.6794 & 0.8489 & 1.0473 & 1.2987 & 1.6759 & 2.0086 & 2.1087 & 2.4033 & 2.6778 \\ 
  60 & 0.6786 & 0.8477 & 1.0455 & 1.2958 & 1.6706 & 2.0003 & 2.0994 & 2.3901 & 2.6603 \\ 
  80 & 0.6776 & 0.8461 & 1.0432 & 1.2922 & 1.6641 & 1.9901 & 2.0878 & 2.3739 & 2.6387 \\ 
  100 & 0.6770 & 0.8452 & 1.0418 & 1.2901 & 1.6602 & 1.9840 & 2.0809 & 2.3642 & 2.6259 \\ 
  1000 & 0.6747 & 0.8420 & 1.0370 & 1.2824 & 1.6464 & 1.9623 & 2.0564 & 2.3301 & 2.5808 \\ 
  $\infty$ & 0.6745 & 0.8416 & 1.0364 & 1.2816 & 1.6449 & 1.9600 & 2.0537 & 2.3263 & 2.5758 \\ 
  \hline
%   \cline{2-10}
  & 50\% & 60\% & 70\% & 80\% & 90\% & 95\% & 96\% & 98\% & 99\%  \\
  \cline{2-10}
 & \multicolumn{9}{ c }{Nível de Confiança $\alpha$} \\
 \cline{2-10}
\end{tabular}
\end{table}



% \newpage
\pagebreak

%\thispagestyle{empty}
\thispagestyle{milktablestyle}

\vspace*{-10mm}
\rput{90}(6.5,-1){
\psscalebox{.75}{
\psset{xunit=1.25cm,yunit=15cm,labelFontSize=\scriptsize}% 
\begin{pspicture}(0,0)(9.25,.25) 
\def\vari{.35}
\pscustom[fillstyle=vlines,hatchwidth=0.1pt,hatchsep=1pt,linestyle=none]{%
\psChiIIDist[linecolor=black, nue=4.5]{6}{9}    
\psline(9,0)(6,0)
}%   
%\psaxes[labels=none,ticks=none,% 
%	arrowscale=1.2]{-}(0,0)(4,0)%
\psaxes[labels=none,ticks=none, arrowscale=1.2]{-}(0,0)(9.25,0)
\psChiIIDist[linecolor=green!60!black,linewidth=1pt,plotpoints=100, nue=4.5]{0.01}{9}    
\rput[c]{0}(0,-.03){\scriptsize\Red{$0$}}% 
\rput[c]{0}(6,-.03){$y_c$}% 
\psline[linecolor=red](0,0)(0,-.01)% 
\psline(6,0)(6,-.01)% 
\rput[l]{0}(5,0.17){\parbox{6.5cm}{Probabilidade $p$}}%
\psline{o-}(7,.03)(6,.15)
\end{pspicture}}%
}

\begin{table}[ht]
\vspace*{-1.75in}\small
\centering
\caption{Valores críticos da distribuição $\chi^2$.}
\begin{sideways}
\begin{tabular}{r||rrrrrrrrrrrrrrrrr}
\hline
 \multirow{2}{*}{gl} & \multicolumn{15}{|c}{Probabilidade de cauda superior $p$} \\
  \cline{2-16}
 & 0.99 & 0.98 & 0.975 & 0.95 & 0.90 & 0.80 & 0.70 & 0.50 & 0.30 & 0.20 & 0.10 & 0.05 & 0.025 & 0.02 & 0.01 \\ 
  \hline
1 & 0.0002 & 0.0006 & 0.0010 & 0.0039 & 0.0158 & 0.0642 & 0.1485 & 0.4549 & 1.0742 & 1.6424 & 2.7055 & 3.8415 & 5.0239 & 5.4119 & 6.6349 \\ 
  2 & 0.0201 & 0.0404 & 0.0506 & 0.1026 & 0.2107 & 0.4463 & 0.7133 & 1.3863 & 2.4079 & 3.2189 & 4.6052 & 5.9915 & 7.3778 & 7.8240 & 9.2103 \\ 
  3 & 0.1148 & 0.1848 & 0.2158 & 0.3518 & 0.5844 & 1.0052 & 1.4237 & 2.3660 & 3.6649 & 4.6416 & 6.2514 & 7.8147 & 9.3484 & 9.8374 & 11.3449 \\ 
  4 & 0.2971 & 0.4294 & 0.4844 & 0.7107 & 1.0636 & 1.6488 & 2.1947 & 3.3567 & 4.8784 & 5.9886 & 7.7794 & 9.4877 & 11.1433 & 11.6678 & 13.2767 \\ 
  5 & 0.5543 & 0.7519 & 0.8312 & 1.1455 & 1.6103 & 2.3425 & 2.9999 & 4.3515 & 6.0644 & 7.2893 & 9.2364 & 11.0705 & 12.8325 & 13.3882 & 15.0863 \\ 
  6 & 0.8721 & 1.1344 & 1.2373 & 1.6354 & 2.2041 & 3.0701 & 3.8276 & 5.3481 & 7.2311 & 8.5581 & 10.6446 & 12.5916 & 14.4494 & 15.0332 & 16.8119 \\ 
  7 & 1.2390 & 1.5643 & 1.6899 & 2.1673 & 2.8331 & 3.8223 & 4.6713 & 6.3458 & 8.3834 & 9.8032 & 12.0170 & 14.0671 & 16.0128 & 16.6224 & 18.4753 \\ 
  8 & 1.6465 & 2.0325 & 2.1797 & 2.7326 & 3.4895 & 4.5936 & 5.5274 & 7.3441 & 9.5245 & 11.0301 & 13.3616 & 15.5073 & 17.5345 & 18.1682 & 20.0902 \\ 
  9 & 2.0879 & 2.5324 & 2.7004 & 3.3251 & 4.1682 & 5.3801 & 6.3933 & 8.3428 & 10.6564 & 12.2421 & 14.6837 & 16.9190 & 19.0228 & 19.6790 & 21.6660 \\ 
  10 & 2.5582 & 3.0591 & 3.2470 & 3.9403 & 4.8652 & 6.1791 & 7.2672 & 9.3418 & 11.7807 & 13.4420 & 15.9872 & 18.3070 & 20.4832 & 21.1608 & 23.2093 \\ 
  11 & 3.0535 & 3.6087 & 3.8157 & 4.5748 & 5.5778 & 6.9887 & 8.1479 & 10.3410 & 12.8987 & 14.6314 & 17.2750 & 19.6751 & 21.9200 & 22.6179 & 24.7250 \\ 
  12 & 3.5706 & 4.1783 & 4.4038 & 5.2260 & 6.3038 & 7.8073 & 9.0343 & 11.3403 & 14.0111 & 15.8120 & 18.5493 & 21.0261 & 23.3367 & 24.0540 & 26.2170 \\ 
  13 & 4.1069 & 4.7654 & 5.0088 & 5.8919 & 7.0415 & 8.6339 & 9.9257 & 12.3398 & 15.1187 & 16.9848 & 19.8119 & 22.3620 & 24.7356 & 25.4715 & 27.6882 \\ 
  14 & 4.6604 & 5.3682 & 5.6287 & 6.5706 & 7.7895 & 9.4673 & 10.8215 & 13.3393 & 16.2221 & 18.1508 & 21.0641 & 23.6848 & 26.1189 & 26.8728 & 29.1412 \\ 
  15 & 5.2293 & 5.9849 & 6.2621 & 7.2609 & 8.5468 & 10.3070 & 11.7212 & 14.3389 & 17.3217 & 19.3107 & 22.3071 & 24.9958 & 27.4884 & 28.2595 & 30.5779 \\ 
  16 & 5.8122 & 6.6142 & 6.9077 & 7.9616 & 9.3122 & 11.1521 & 12.6243 & 15.3385 & 18.4179 & 20.4651 & 23.5418 & 26.2962 & 28.8454 & 29.6332 & 31.9999 \\ 
  17 & 6.4078 & 7.2550 & 7.5642 & 8.6718 & 10.0852 & 12.0023 & 13.5307 & 16.3382 & 19.5110 & 21.6146 & 24.7690 & 27.5871 & 30.1910 & 30.9950 & 33.4087 \\ 
  18 & 7.0149 & 7.9062 & 8.2307 & 9.3905 & 10.8649 & 12.8570 & 14.4399 & 17.3379 & 20.6014 & 22.7595 & 25.9894 & 28.8693 & 31.5264 & 32.3462 & 34.8053 \\ 
  19 & 7.6327 & 8.5670 & 8.9065 & 10.1170 & 11.6509 & 13.7158 & 15.3517 & 18.3377 & 21.6891 & 23.9004 & 27.2036 & 30.1435 & 32.8523 & 33.6874 & 36.1909 \\ 
  20 & 8.2604 & 9.2367 & 9.5908 & 10.8508 & 12.4426 & 14.5784 & 16.2659 & 19.3374 & 22.7745 & 25.0375 & 28.4120 & 31.4104 & 34.1696 & 35.0196 & 37.5662 \\ 
  21 & 8.8972 & 9.9146 & 10.2829 & 11.5913 & 13.2396 & 15.4446 & 17.1823 & 20.3372 & 23.8578 & 26.1711 & 29.6151 & 32.6706 & 35.4789 & 36.3434 & 38.9322 \\ 
  22 & 9.5425 & 10.6000 & 10.9823 & 12.3380 & 14.0415 & 16.3140 & 18.1007 & 21.3370 & 24.9390 & 27.3015 & 30.8133 & 33.9244 & 36.7807 & 37.6595 & 40.2894 \\ 
  23 & 10.1957 & 11.2926 & 11.6886 & 13.0905 & 14.8480 & 17.1865 & 19.0211 & 22.3369 & 26.0184 & 28.4288 & 32.0069 & 35.1725 & 38.0756 & 38.9683 & 41.6384 \\ 
  24 & 10.8564 & 11.9918 & 12.4012 & 13.8484 & 15.6587 & 18.0618 & 19.9432 & 23.3367 & 27.0960 & 29.5533 & 33.1962 & 36.4150 & 39.3641 & 40.2704 & 42.9798 \\ 
  25 & 11.5240 & 12.6973 & 13.1197 & 14.6114 & 16.4734 & 18.9398 & 20.8670 & 24.3366 & 28.1719 & 30.6752 & 34.3816 & 37.6525 & 40.6465 & 41.5661 & 44.3141 \\ 
  26 & 12.1981 & 13.4086 & 13.8439 & 15.3792 & 17.2919 & 19.8202 & 21.7924 & 25.3365 & 29.2463 & 31.7946 & 35.5632 & 38.8851 & 41.9232 & 42.8558 & 45.6417 \\ 
  27 & 12.8785 & 14.1254 & 14.5734 & 16.1514 & 18.1139 & 20.7030 & 22.7192 & 26.3363 & 30.3193 & 32.9117 & 36.7412 & 40.1133 & 43.1945 & 44.1400 & 46.9629 \\ 
  28 & 13.5647 & 14.8475 & 15.3079 & 16.9279 & 18.9392 & 21.5880 & 23.6475 & 27.3362 & 31.3909 & 34.0266 & 37.9159 & 41.3371 & 44.4608 & 45.4188 & 48.2782 \\ 
  29 & 14.2565 & 15.5745 & 16.0471 & 17.7084 & 19.7677 & 22.4751 & 24.5770 & 28.3361 & 32.4612 & 35.1394 & 39.0875 & 42.5570 & 45.7223 & 46.6927 & 49.5879 \\ 
  30 & 14.9535 & 16.3062 & 16.7908 & 18.4927 & 20.5992 & 23.3641 & 25.5078 & 29.3360 & 33.5302 & 36.2502 & 40.2560 & 43.7730 & 46.9792 & 47.9618 & 50.8922 \\ 
   \hline
  %%%%%%%%%%%%%%%%%%%%%%%%%55
  \hline
%   \cline{2-10}
 & 0.99 & 0.98 & 0.975 & 0.95 & 0.9 & 0.8 & 0.7 & 0.5 & 0.3 & 0.2 & 0.1 & 0.05 & 0.025 & 0.02 & 0.01 \\ 
  \cline{2-16}
\end{tabular}
\end{sideways}
\end{table}



% \newpage
\pagebreak

%\thispagestyle{empty}
\thispagestyle{milktablestyle}

\vspace*{-10mm}
\rput{90}(6.5,-1){
\psscalebox{.75}{
\psset{xunit=1.25cm,yunit=15cm,labelFontSize=\scriptsize}% 
\begin{pspicture}(0,0)(9.25,.25) 
\def\vari{.35}
\pscustom[fillstyle=vlines,hatchwidth=0.1pt,hatchsep=1pt,linestyle=none]{%
\psFDist[linecolor=black, nue=3,mue=10]{4}{9}    
\psline(9,0)(4,0)
}%   
%\psaxes[labels=none,ticks=none,% 
%	arrowscale=1.2]{-}(0,0)(4,0)%
\psaxes[labels=none,ticks=none, arrowscale=1.2]{-}(0,0)(9.25,0)
\psFDist[linecolor=green!60!black,linewidth=1pt,plotpoints=100, nue=3,mue=10]{0.01}{9}    
\rput[c]{0}(0,-.03){\scriptsize\Red{$0$}}% 
\rput[c]{0}(4,-.03){$f_c$}% 
\psline[linecolor=red](0,0)(0,-.01)% 
\psline(4,0)(4,-.01)% 
\rput[l]{0}(5,0.17){\parbox{6.5cm}{Probabilidade $p$}}%
\psline{o-}(4.5,.02)(6,.15)
\end{pspicture}}%
}

\begin{table}[ht]
\vspace*{-1.75in}\scriptsize
\def\cw{5mm}
\centering
\caption{Valores críticos da distribuição $F$.}
\begin{sideways}
\begin{tabular}{r||rrrrrrrrrrrrrrrrrrrrrr}
\hline
 \multirow{2}{*}{gl} & \multicolumn{22}{c}{Probabilidade de cauda superior $p=5\%$. $P(f > f_c) = 5\%$} \\
  \cline{2-21}
 & 1 & 2 & 3 & 4 & 5 & 6 & 7 & 8 & 9 & 10 & 12 & 14 & 15 & 16 & 18 & 20 & 24 & 30 & 40 & 60 & 120 & $\infty$ \\ 
  \hline
1 & 161.4 & 199.5 & 215.7 & 224.6 & 230.2 & 234.0 & 236.8 & 238.9 & 240.5 & 241.9 & 243.9 & 245.4 & 246.0 & 246.5 & 247.3 & 248.0 & 249.0 & 250.1 & 251.1 & 252.2 & 253.2 & 254.3 \\ 
  2 & 18.51 & 19.00 & 19.16 & 19.25 & 19.30 & 19.33 & 19.35 & 19.37 & 19.38 & 19.40 & 19.41 & 19.42 & 19.43 & 19.43 & 19.44 & 19.45 & 19.45 & 19.46 & 19.47 & 19.48 & 19.49 & 19.50 \\ 
  3 & 10.13 & 9.55 & 9.28 & 9.12 & 9.01 & 8.94 & 8.89 & 8.85 & 8.81 & 8.79 & 8.74 & 8.71 & 8.70 & 8.69 & 8.67 & 8.66 & 8.64 & 8.62 & 8.59 & 8.57 & 8.55 & 8.53 \\ 
  4 & 7.71 & 6.94 & 6.59 & 6.39 & 6.26 & 6.16 & 6.09 & 6.04 & 6.00 & 5.96 & 5.91 & 5.87 & 5.86 & 5.84 & 5.82 & 5.80 & 5.77 & 5.75 & 5.72 & 5.69 & 5.66 & 5.63 \\ 
  5 & 6.61 & 5.79 & 5.41 & 5.19 & 5.05 & 4.95 & 4.88 & 4.82 & 4.77 & 4.74 & 4.68 & 4.64 & 4.62 & 4.60 & 4.58 & 4.56 & 4.53 & 4.50 & 4.46 & 4.43 & 4.40 & 4.37 \\ 
  6 & 5.99 & 5.14 & 4.76 & 4.53 & 4.39 & 4.28 & 4.21 & 4.15 & 4.10 & 4.06 & 4.00 & 3.96 & 3.94 & 3.92 & 3.90 & 3.87 & 3.84 & 3.81 & 3.77 & 3.74 & 3.70 & 3.67 \\ 
  7 & 5.59 & 4.74 & 4.35 & 4.12 & 3.97 & 3.87 & 3.79 & 3.73 & 3.68 & 3.64 & 3.57 & 3.53 & 3.51 & 3.49 & 3.47 & 3.44 & 3.41 & 3.38 & 3.34 & 3.30 & 3.27 & 3.23 \\ 
  8 & 5.32 & 4.46 & 4.07 & 3.84 & 3.69 & 3.58 & 3.50 & 3.44 & 3.39 & 3.35 & 3.28 & 3.24 & 3.22 & 3.20 & 3.17 & 3.15 & 3.12 & 3.08 & 3.04 & 3.01 & 2.97 & 2.93 \\ 
  9 & 5.12 & 4.26 & 3.86 & 3.63 & 3.48 & 3.37 & 3.29 & 3.23 & 3.18 & 3.14 & 3.07 & 3.03 & 3.01 & 2.99 & 2.96 & 2.94 & 2.90 & 2.86 & 2.83 & 2.79 & 2.75 & 2.71 \\ 
  10 & 4.96 & 4.10 & 3.71 & 3.48 & 3.33 & 3.22 & 3.14 & 3.07 & 3.02 & 2.98 & 2.91 & 2.86 & 2.85 & 2.83 & 2.80 & 2.77 & 2.74 & 2.70 & 2.66 & 2.62 & 2.58 & 2.54 \\ 
  12 & 4.75 & 3.89 & 3.49 & 3.26 & 3.11 & 3.00 & 2.91 & 2.85 & 2.80 & 2.75 & 2.69 & 2.64 & 2.62 & 2.60 & 2.57 & 2.54 & 2.51 & 2.47 & 2.43 & 2.38 & 2.34 & 2.30 \\ 
  14 & 4.60 & 3.74 & 3.34 & 3.11 & 2.96 & 2.85 & 2.76 & 2.70 & 2.65 & 2.60 & 2.53 & 2.48 & 2.46 & 2.44 & 2.41 & 2.39 & 2.35 & 2.31 & 2.27 & 2.22 & 2.18 & 2.13 \\ 
  15 & 4.54 & 3.68 & 3.29 & 3.06 & 2.90 & 2.79 & 2.71 & 2.64 & 2.59 & 2.54 & 2.48 & 2.42 & 2.40 & 2.38 & 2.35 & 2.33 & 2.29 & 2.25 & 2.20 & 2.16 & 2.11 & 2.07 \\ 
  16 & 4.49 & 3.63 & 3.24 & 3.01 & 2.85 & 2.74 & 2.66 & 2.59 & 2.54 & 2.49 & 2.42 & 2.37 & 2.35 & 2.33 & 2.30 & 2.28 & 2.24 & 2.19 & 2.15 & 2.11 & 2.06 & 2.01 \\ 
  18 & 4.41 & 3.55 & 3.16 & 2.93 & 2.77 & 2.66 & 2.58 & 2.51 & 2.46 & 2.41 & 2.34 & 2.29 & 2.27 & 2.25 & 2.22 & 2.19 & 2.15 & 2.11 & 2.06 & 2.02 & 1.97 & 1.92 \\ 
  20 & 4.35 & 3.49 & 3.10 & 2.87 & 2.71 & 2.60 & 2.51 & 2.45 & 2.39 & 2.35 & 2.28 & 2.22 & 2.20 & 2.18 & 2.15 & 2.12 & 2.08 & 2.04 & 1.99 & 1.95 & 1.90 & 1.84 \\ 
  24 & 4.26 & 3.40 & 3.01 & 2.78 & 2.62 & 2.51 & 2.42 & 2.36 & 2.30 & 2.25 & 2.18 & 2.13 & 2.11 & 2.09 & 2.05 & 2.03 & 1.98 & 1.94 & 1.89 & 1.84 & 1.79 & 1.73 \\ 
  30 & 4.17 & 3.32 & 2.92 & 2.69 & 2.53 & 2.42 & 2.33 & 2.27 & 2.21 & 2.16 & 2.09 & 2.04 & 2.01 & 1.99 & 1.96 & 1.93 & 1.89 & 1.84 & 1.79 & 1.74 & 1.68 & 1.62 \\ 
  40 & 4.08 & 3.23 & 2.84 & 2.61 & 2.45 & 2.34 & 2.25 & 2.18 & 2.12 & 2.08 & 2.00 & 1.95 & 1.92 & 1.90 & 1.87 & 1.84 & 1.79 & 1.74 & 1.69 & 1.64 & 1.58 & 1.51 \\ 
  60 & 4.00 & 3.15 & 2.76 & 2.53 & 2.37 & 2.25 & 2.17 & 2.10 & 2.04 & 1.99 & 1.92 & 1.86 & 1.84 & 1.82 & 1.78 & 1.75 & 1.70 & 1.65 & 1.59 & 1.53 & 1.47 & 1.39 \\ 
  120 & 3.92 & 3.07 & 2.68 & 2.45 & 2.29 & 2.18 & 2.09 & 2.02 & 1.96 & 1.91 & 1.83 & 1.78 & 1.75 & 1.73 & 1.69 & 1.66 & 1.61 & 1.55 & 1.50 & 1.43 & 1.35 & 1.26 \\ 
  $\infty$ & 3.84 & 3.00 & 2.61 & 2.37 & 2.21 & 2.10 & 2.01 & 1.94 & 1.88 & 1.83 & 1.75 & 1.69 & 1.67 & 1.64 & 1.60 & 1.57 & 1.52 & 1.46 & 1.40 & 1.32 & 1.22 & 1.03 \\ 
   \hline
   \hline
 & 1 & 2 & 3 & 4 & 5 & 6 & 7 & 8 & 9 & 10 & 12 & 14 & 15 & 16 & 18 & 20 & 24 & 30 & 40 & 60 & 120 & $\infty$ \\
 \hline
\end{tabular}
\end{sideways}
\end{table}



\end{document}

% Prova
